For $0\le s<Q$ put

\begin{equation}
 B_s(x)=\binom{Qx-s-1}{hQ}.
\end{equation}
Pair the nonmultiples of $Q$ in the numerator window with the corresponding
nonmultiples in $\{1,\ldots,hQ\}$. Each quotient is a $p$-adic unit whose
reduction does not depend on $x$. So the product formula first shows $B_s(x)\in\mathbb Z_{(p)}[x]$.
Partition the $hQ$ numerator factors
according to their residue modulo $Q$. Every residue occurs $h$ times.
More explicitly, pair a numerator factor $Qx-c$ with a denominator factor
$c'$ having the same residue modulo $Q$. Then $c\equiv c'\pmod p$ and
$(Qx-c)/c'\equiv-1\pmod p$. The $h$ factors with $Q\mid c$ give the
remaining product $\binom{x-1}{h}$ after the powers of $Q$ are canceled.
After removing the fixed powers of $p$ contributed by $(hQ)!$, all residue
classes except the one with $t\equiv-s-1\pmod Q$ contribute $p$-adic units
independent of $x$ modulo $p$. The exceptional class contributes exactly
$h$ linear factors in $x$. Thus the reduced polynomial has degree at most
$h$. Lucas' formula gives its values at $x=1,\ldots,h$. At $x=0$ the
negative-index identity gives the same value $(-1)^h$, because
$\binom{-1}{h}=(-1)^h$; on the left $Q=p^k$ is odd since $p>2n$.
The remaining sign is $(-1)^{h(Q-1)}=1$. The $h+1$ values therefore prove the coefficientwise congruence
is the unsigned identity
\begin{equation}
 \binom{Qx-s-1}{hQ}\equiv\binom{x-1}{h}\pmod p.       \label{eq:20}
\end{equation}

Write $i=Qi_0+s_i$ and $j=Qj_0+s_j$, with residues in
$\{0,\ldots,Q-1\}$. Let $P^0_{i_0j_0}$, $E^0=(E^0_{i_0j_0})$ and
$C^0=(C^0_{i_0j_0})$ be the polynomial parts $P_{ij}$ and the matrices
$E(a)$ and $C(a)$ of Section~\ref{sec:determinant} with $(m,g)=(n,h)$.
For a nonzero offset $d$ we use the Taylor series
\[
 U_d(x)=\frac{Q}{Qx+d}=\frac{Q/d}{1+(Q/d)x}.
\]
Since $|d|<2nQ<pQ$, its coefficients are $p$-integral. Put
$d_i=w(-i)/(j-i)$ and $d_j=-w(-j)/(j-i)$. The scaled polynomial part
satisfies the exact identity
\begin{equation}
 Q^2P_{ij}(Qx-s)=B_s(x)^3U_{i-s}(x)U_{j-s}(x)
 -(Qd_i)U_{i-s}(x)-(Qd_j)U_{j-s}(x).                  \label{eq:20a}
\end{equation}
Suppose one offset is zero. The two residue values then agree. So the
numerator in the corresponding $1/x$ term has zero constant coefficient.
Division by $x$ therefore remains integral. Now reduce \eqref{eq:20a}. A nonzero
offset contributes only when it is a multiple of $Q$. Both offsets
contribute precisely when $s_i=s_j=s$. In that case \eqref{eq:20} gives the seed
polynomial part $P^0_{i_0j_0}$. This includes the finite-part division at a
zero offset. After two derivatives we obtain, coefficientwise,
\begin{equation}
 Q^4P_{ij}''(Qx-s)\equiv
 {\bf1}_{s_i=s_j=s}(P^0_{i_0j_0})''(x)\pmod p.             \label{eq:21}
\end{equation}

The scaled values in \eqref{eq:21} have integral Newton coefficients. So Newton
interpolation and
\[
 \sum_{r\ge1}z^r\binom{r-1}{t}=
 \frac{z^{t+1}}{(1-z)^{t+1}}
\]
sum \eqref{eq:21} coefficientwise. Split $r=Qx-s$. A matching residue then
contributes the factor $a^s$ and the denominator $1-a^{-Q}$. This
denominator is a $p$-unit because $p\nmid a(a-1)$, so
$a^Q\equiv a\not\equiv1\pmod p$. The
zero-offset finite part is included before this splitting. Because $h<p$,
the $p$-parts of the factors in $(hQ)!$ cancel coefficientwise against the
corresponding numerator blocks. For the harmonic term,
\[
 Q^3H_{Qu+s}(a)\equiv a^sH_u(a)\pmod p,
\]
because only the multiples $t=Qv$ survive and $a^Q=a$. We combine the
polynomial part, the harmonic term, and the $T$ coefficient. This gives the
all-$k$ congruences
\begin{equation}
 Q^4E_{ij}\equiv {\bf1}_{s_i=s_j}a^sE^0_{i_0j_0},\qquad
 QC_{ij}\equiv {\bf1}_{s_i=s_j}a^sC^0_{i_0j_0}\pmod p.       \label{eq:22}
\end{equation}

Define
\[
 {\mathcal P}_k(X)=Q^{4m}G_{m,a}(X/Q^3).
\]
Group rows and columns by their residues modulo $Q$. The block permutation
signs cancel. The $Q$ diagonal blocks are
$a^s(E^0+XC^0)$. Their scalar is
$a^{nQ(Q-1)/2}=1$ in $\mathbb F_p$, since $n$ is even and
$a^{Q-1}=1$. Thus, for every $k$,
\begin{equation}
 {\mathcal P}_k(X)\in\mathbb Z_{(p)}[X],\qquad
 {\mathcal P}_k(X)\equiv G_{n,a}(X)^Q\pmod p.                  \label{eq:23}
\end{equation}
Let $\xi$ be fixed and rational. Choose $k$ with
$3k+v_p(\xi)\ge1$ (with $v_p(0)=+\infty$). Evaluate at
$X=Q^3\xi$. For all sufficiently large $k$ this gives the unscaled identity
\[
 Q^{4m}G_{m,a}(\xi)={\mathcal P}_k(Q^3\xi)
 \equiv G_{n,a}(Q^3\xi)^Q
 \equiv G_{n,a}(0)^Q\not\equiv0\pmod p.
\]
The seed choice in \eqref{eq:19} makes $G_{n,a}(0)$ a $p$-unit. Since the
left side would be zero if $G_{m,a}(\xi)=0$, this unit congruence gives the
required contradiction. The argument uses the displayed congruence before
any cancellation of the factor $Q^{4m}$.
